\begin{figure}[!htbp]
  \centering
  \resizebox{0.98\linewidth}{!}{%
  \begin{tikzpicture}[
    x=1cm, y=1cm,
    >=Stealth,
    every node/.style={font=\rmfamily\small, align=center, text=black},
    stage/.style={
      rectangle, rounded corners=6pt,
      draw=goldTone!42!black, fill=black!1, line width=0.65pt
    },
    extStage/.style={
      rectangle, rounded corners=6pt,
      draw=black!32, dashed, fill=black!2, line width=0.62pt
    },
    stageLabel/.style={font=\bfseries\scriptsize, text=goldTone!64!black, anchor=west},
    extLabel/.style={font=\bfseries\scriptsize, text=black!58, anchor=west},
    profCard/.style={
      rectangle, rounded corners=5pt,
      draw=greenTone!58!black, fill=greenFill, line width=0.72pt
    },
    profCardSoft/.style={
      rectangle, rounded corners=5pt,
      draw=greenTone!50!black, fill=white, line width=0.68pt
    },
    goldCard/.style={
      rectangle, rounded corners=5pt,
      draw=goldTone!58!black, fill=goldFill, line width=0.72pt
    },
    cardTitleGold/.style={font=\bfseries\scriptsize, text=goldTone!58!black},
    cardTitleGreen/.style={font=\bfseries\scriptsize, text=greenTone!54!black},
    note/.style={font=\scriptsize, text=black!56},
    noteSmall/.style={font=\tiny, text=black!58},
    arr/.style={-Stealth, draw=goldTone!58!black, line width=0.80pt}
  ]
    \definecolor{goldTone}{RGB}{178,132,37}
    \definecolor{goldFill}{RGB}{251,240,211}
    \definecolor{greenTone}{RGB}{66,134,108}
    \definecolor{greenFill}{RGB}{229,243,236}
    \definecolor{blueTone}{RGB}{62,103,167}

    % ==========================================
    % 左 panel：BA-k 与 AutoK 主线分配机制
    % ==========================================
    \node[stage, minimum width=11.20cm, minimum height=5.30cm] (mainStage) at (5.75,2.70) {};
    \node[stageLabel] at (0.30,5.15) {（a）BA-$k$ 与 AutoK 主线分配机制};

    % --- Card 1：校准产物输入（绿，与图 3-3 画像同色）
    \node[profCardSoft, minimum width=2.10cm, minimum height=3.10cm] (input) at (1.45,2.50) {};
    \node[cardTitleGreen] at (1.45,3.85) {校准产物输入};
    \node[font=\scriptsize] at (1.45,3.30) {$\sigma_K^{(l)}$};
    \node[note] at (1.45,2.85) {K-path};
    \node[note] at (1.45,2.50) {per-channel};
    \node[note] at (1.45,2.05) {scale 范围};
    \node[noteSmall] at (1.45,1.30) {离线副产物};

    % --- Card 2：画像聚合 𝒮（绿，主体）
    \node[profCard, minimum width=2.50cm, minimum height=3.10cm] (profile) at (4.40,2.50) {};
    \node[cardTitleGreen] at (4.40,3.85) {层级敏感度画像 $\mathcal S$};
    \node[font=\scriptsize] at (4.40,3.36) {$s^{(l)}=\mathrm{Agg}_\alpha(\sigma_K^{(l)})$};
    % σ_K 画像 bar chart（8 层示意）
    \begin{scope}[shift={(3.55,1.55)}]
      \foreach \x/\h in {0/0.42, 1/0.66, 2/0.30, 3/0.88, 4/0.55, 5/0.26, 6/0.78, 7/0.40} {
        \filldraw[fill=greenFill!40!greenTone, draw=greenTone!72!black, line width=0.32pt]
          ({0.20*\x},0) rectangle ++(0.16,\h);
      }
    \end{scope}
    \node[noteSmall] at (4.40,1.30) {层 $1\,\dots\,L$};
    \node[note] at (4.40,1.95) {$\alpha=\max$};

    % --- Card 3：Γ(k) + AutoK（金色，视觉焦点）
    \node[goldCard, minimum width=3.10cm, minimum height=3.10cm] (gamma) at (8.05,2.50) {};
    \node[cardTitleGold] at (8.05,3.85) {覆盖度 $\Gamma(k)$ 与 AutoK};
    % Γ(k) 阶梯曲线 + ρ + k* 反推
    \begin{scope}[shift={(6.85,1.40)}]
      % 坐标轴
      \draw[black!50, line width=0.40pt] (0,0) -- (2.30,0);
      \draw[black!50, line width=0.40pt] (0,0) -- (0,1.20);
      % 阶梯曲线 Γ(k)：单调上升至 1
      \draw[goldTone, line width=1.0pt]
        (0,0) -- (0.20,0.32) -- (0.40,0.32) -- (0.40,0.55)
        -- (0.60,0.55) -- (0.60,0.72) -- (0.80,0.72)
        -- (0.80,0.88) -- (1.00,0.88) -- (1.00,0.96)
        -- (1.20,0.96) -- (1.20,1.00) -- (2.30,1.00);
      % ρ 阈值水平虚线
      \draw[blueTone, dashed, line width=0.55pt] (-0.05,0.85) -- (2.30,0.85);
      \node[font=\tiny, text=blueTone!72!black, anchor=west] at (2.05,0.93) {$\rho$};
      % k* 反推垂直虚线
      \draw[goldTone!75!black, dashed, line width=0.55pt] (0.85,0) -- (0.85,0.88);
      \node[font=\tiny, text=goldTone!72!black, anchor=north] at (0.85,-0.05) {$k^\star\!(\rho)$};
      % 轴标签
      \node[font=\tiny, text=black!58, anchor=east] at (-0.04,1.10) {$1$};
      \node[font=\tiny, text=black!58, anchor=north east] at (-0.04,0.85) {};
      \node[font=\tiny, text=black!58, anchor=east] at (-0.04,0.40) {$\Gamma$};
      \node[font=\tiny, text=black!58, anchor=north] at (2.30,-0.05) {$k$};
    \end{scope}
    \node[noteSmall] at (8.05,1.10) {$k^\star\!=\!\min\{k:\Gamma(k)\!\ge\!\rho\}$};

    % --- Card 4：两级阶梯位宽 b^(l)（金色）
    \node[goldCard, minimum width=2.40cm, minimum height=3.10cm] (out) at (10.85,2.50) {};
    \node[cardTitleGold] at (10.85,3.85) {两级位宽 $b^{(l)}$};
    % 8 层位宽阶梯：未保护层 b_lo=4（短），保护层 b_hi=8（长）
    \begin{scope}[shift={(10.05,1.40)}]
      % 低 bit 层（4 bit）
      \foreach \x in {0, 1, 3, 5, 6} {
        \filldraw[fill=white, draw=goldTone!50, line width=0.42pt]
          ({0.20*\x},0) rectangle ++(0.16,0.36);
      }
      % 高 bit 层（8 bit）— 第 2、4、7 层是保护集合 P_k
      \foreach \x in {2, 4, 7} {
        \filldraw[fill=goldTone!45, draw=goldTone!82!black, line width=0.48pt]
          ({0.20*\x},0) rectangle ++(0.16,0.78);
      }
      % 4/8 标签（右侧）
      \node[font=\tiny, text=black!56, anchor=west] at (1.65,0.78) {$8$};
      \node[font=\tiny, text=black!56, anchor=west] at (1.65,0.36) {$4$};
    \end{scope}
    \node[noteSmall] at (10.85,1.10) {$b\!\in\!\{4,8\}$};
    \node[noteSmall] at (10.85,0.78) {$\mathcal P_k=\mathrm{TopK}(\mathcal S,k)$};

    % --- 主线箭头
    \draw[arr] (input.east) -- node[above, font=\tiny, text=goldTone!64!black] {聚合} (profile.west);
    \draw[arr] (profile.east) -- node[above, font=\tiny, text=goldTone!64!black] {归一/累积} (gamma.west);
    \draw[arr] (gamma.east) -- node[above, font=\tiny, text=goldTone!64!black] {$\mathrm{TopK}$} (out.west);

    % ==========================================
    % 右 panel：K/V 角色接口预留（虚线小框）
    % ==========================================
    \node[extStage, minimum width=2.85cm, minimum height=5.30cm] (extP) at (12.95,2.70) {};
    \node[extLabel] at (11.55,5.15) {（b）K/V 角色接口};

    \node[font=\bfseries\scriptsize, text=black!58] at (12.95,4.45) {扩展预留\,（\S\,\ref{subsec:ch3-budget-kv}）};

    \node[font=\scriptsize] at (12.95,3.85) {$\mathbf b^{(l)}\!=\!(b_K^{(l)},b_V^{(l)})$};
    \node[note] at (12.95,3.30) {独立 $k_K$、$k_V$};

    \draw[black!22, line width=0.50pt] (11.60,2.92) -- (14.30,2.92);

    \node[note, text=black!58] at (12.95,2.62) {主线设置};
    \node[font=\scriptsize, text=black!62] at (12.95,2.20) {$k_K\!=\!k_V\!=\!k$};
    \node[note, text=black!54, font=\scriptsize\itshape] at (12.95,1.78) {退化为左侧 BA-$k$};

    \draw[black!22, line width=0.50pt] (11.60,1.45) -- (14.30,1.45);

    \node[note, text=black!54, font=\scriptsize\itshape] at (12.95,1.12) {$k_K\!\neq\!k_V$};
    \node[note, text=black!54, font=\scriptsize\itshape] at (12.95,0.78) {留待后续工作};

  \end{tikzpicture}%
  }
  \caption{行为敏感度画像到预算分配的两级机制。\textbf{(a) 主线}：从校准产物 \(\sigma_K^{(l)}\)（K 侧 per-channel scale 范围）出发，按 \(\alpha=\max\) 聚合得到层级敏感度画像 \(\mathcal S=\{s^{(l)}\}_{l=1}^{L}\)；归一化与累积后形成覆盖度阶梯函数 \(\Gamma(k)\)，AutoK 由覆盖率阈值 \(\rho\) 反推满足 \(\Gamma(k)\!\ge\!\rho\) 的最小保护层数 \(k^\star\)（BA-\(k\) 则按人工指定整数 \(k\)）；以 \(\mathrm{TopK}\) 选出保护集合 \(\mathcal P_k\) 后，按两级映射写出位宽向量 \(b^{(l)}\!\in\!\{4,8\}\)，保护层取 \(b_\mathrm{hi}\!=\!8\)，其余取 \(b_\mathrm{lo}\!=\!4\)。\textbf{(b) 接口预留}：第~\ref{subsec:ch3-budget-kv}~节定义角色感知接口 \(\mathbf b^{(l)}=(b_K^{(l)},b_V^{(l)})\)；本文主线固定 \(k_K\!=\!k_V\!=\!k\) 退化为 BA-\(k\)，\(k_K\!\neq\!k_V\) 的实证比较留待后续工作。}
  \label{fig:ch3-allocation-flow}
\end{figure}
